\section{Cohomology of the completed strata}\label{h3:sec:analytic}

We compute the two completed strata together with their maps from
constants.  These maps retain the class from the rank-two Tate frame
through the changes of coefficients later in the proof.
The Kummer and translation calculations of
Barthel--Mann--Ray--Schlank--Senger--Weinstein--Zhou
\cite[Sections~3.3--3.6]{BMR} provide the methodological setting.
We prove the relative estimates, support maps, and limit arguments
needed to retain the particular constants maps in this setting.

We use the relative extension sheaf $N_m$ and integral-frame stacks
of \Cref{h3:sec:towers}. In this section we abbreviate
\[
 E_m=\mathcal E_m=(N_m^{\,3-m})_{\mathrm{ind}},\qquad
 Y_m=\mathcal Y_m\simeq
 [E_m/(\GL_{3-m}(\Z_p)\times P_m)].
\]
Here $N_m$ is the sheafified relative $H^1$ of $V_m^\vee$ defined
before \Cref{h3:geom:independent-extensions}, and the stack equivalence
is \Cref{h3:prop:integral-frames}. Independence is over $\Q_p$.
In particular $E_2=N_2^*$, whereas
$E_1$ is the space of independent pairs in $N_1$.  All coefficient maps
below are maps on these stacks or on their marking torsors.  They
therefore commute with the actions and the changes of marking in
\Cref{h3:prop:integral-frames}.

\subsection{Relative coefficients and the affine calculation}

Fix an algebraically closed perfectoid untilt $C$ of $C^\flat$ and an
embedding $k\subset C^\flat$.  If $S=\operatorname{Spa}(R,R^+)$ is an
affinoid perfectoid space over $C^\flat$, a relative affine space means
the diamond of the affine space over its induced untilt.  Its
characteristic-$p$ structure sheaf is denoted by $\mathcal O^\flat$.
The notation $R(-1)$ retains the Kummer orientation line; choosing
compatible roots of unity in $C$ identifies it with $R$.  This choice
does not make the arithmetic action on this line trivial.

For the ordinary stratum over this relative base, we will prove that
the structure map
$Y_1\to B(\GL_2(\Z_p)\times P_1)$ induces an equivalence
\[
 R\Gamma_{\hthreects}(\GL_2(\Z_p)\times P_1,R)
   \longrightarrow R\Gamma(Y_1,\mathcal O^\flat).
\]
For coheight one, we first compute $N_2^*$: its cohomology consists
of constants $R$ in degree zero and a free $R$-line in degree three,
whose scalar action must be retained.  We will obtain both descriptions
by computing support cohomology in affine spaces and then taking
translation and frame-group cohomology.  Descent to $k$ will turn these
relative calculations into the completed coefficient comparisons,
including the finite connected-marking stages needed for coheight one.

Our relative assertions also include a profinite parameter space $T$.
On coefficients this replaces $R$ by $C_{\hthreects}(T,R)$.  Products have
their product topology, and countable inverse limits are derived until
their higher derived limits have been shown to vanish.  These
conventions are particularly relevant for the distribution modules
introduced below.
Finiteness means finite-dimensional cohomology over $C^\flat$ at the
geometric base $R=C^\flat$ and $T=*$; after descent, the field is $k$.
For an arbitrary relative base $R$, the corresponding answer is obtained
by scalar extension from $C^\flat$.
The corresponding parameter assertion for a computing complex
$K^\bullet$ is the natural identification
\[
 H^a(C_{\hthreects}(T,K^\bullet))
       =C_{\hthreects}(T,H^a(K^\bullet)).
\]
Thus, after descent to $k$, a finite-dimensional answer gives a
finite free module over $C_{\hthreects}(T,k)$.  It need not be
finite-dimensional over $k$ when $T$ is infinite.  The proofs
below establish this identity for the actual computing complexes.

\begin{lemma}[The relative additive presentation]\label{h3:ana:relative-bc}
Let $F/\Q_p$ be the unramified extension of degree $m$, embedded in $C$.
There is an identification over $C^\flat$
\[
 N_m\simeq \mathbb G_{a,C}^{\diamond}/\underline F.
\]
It is compatible with base change in $S$ and with scalar multiplication
by $\Q_p^\times$.  In particular
\[
 N_m^*\simeq
 [(\mathbb A^1_C\setminus\underline F)^{\diamond}/\underline F].
\]
\end{lemma}
\begin{proof}
On the Fargues--Fontaine curve with coefficient field $F$, the divisor
of the chosen untilt gives
\[
 0\longrightarrow\mathcal O_F(-1)\longrightarrow\mathcal O_F
   \longrightarrow i_*\mathcal O_{S^\sharp}\longrightarrow0.
\]
The relative section sheaf of $\mathcal O_F$ is $\underline F$,
its relative first cohomology vanishes, and
$\mathcal O_F(-1)$ has no relative sections.  These are the
slope-zero and negative-slope cases of the relative bundle
calculation in \Cref{h3:geom:curve-inputs} and \cite{FarguesFontaine,KedlayaLiu}. Pushforward along the
unramified coefficient extension identifies
$\mathcal O_F(-1)$ with $V_m^\vee$: its isocrystal has rank $m$,
degree $-1$, and the cyclic Frobenius of slope $-1/m$.
The cohomology sequence is consequently the displayed quotient.
All its arrows are relative arrows of sheaves, so the presentation
holds over $S$, not only over geometric fields.  The inverse image of
the zero section is precisely $\underline F$.
\end{proof}

Here is the estimate used in the nonproper calculation.  It is useful
to give it explicitly, since affinoid perfectoid acyclicity alone does
not compute the cohomology of an untilted affine line.

\begin{lemma}[Restriction on Kummer annuli]\label{h3:ana:annulus-contraction}
Fix a topological generator of $\Z_p(1)$, let $\epsilon\in C^\flat$
be its compatible system of roots, and put $\theta=|\epsilon-1|<1$.
Consider an untilted closed annulus
$A[r,b]=\{r\leq |z|\leq b\}$ and a smaller annulus
$A[r',b']$ satisfying
\[
 r'\geq r/\theta,\qquad b'\leq\theta b.
\]
In the Kummer cochain models, restriction is homotopic to its
constant-monomial part.  That part is
$R\oplus R(-1)[-1]$.  The homotopy on the nonconstant part is
bounded by one in these two Gauss norms.  The assertion holds
uniformly after $R$ is replaced by $C_{\hthreects}(T,R)$.
\end{lemma}
\begin{proof}
Adjoin all $p$-power roots of $z$.  The resulting annulus is
perfectoid, and its $\Z_p(1)$-torsor computes the cohomology by
\[
 [\,M\xrightarrow{\gamma-1}M\,].
\]
Here $M$ is the tilted Laurent algebra, completed for the two Gauss
norms.  Its elements have unique expansions
$\sum_{u\in\Z[1/p]}a_u z^u$; the weighted coefficients form a
$c_0$ family.  In this statement $c_0$ concerns all the fractional
exponents: for each positive bound only finitely many weighted
coefficients exceed that bound.  It is stronger than a condition
only as $|u|$ tends to infinity.

The differential on the $u$-th monomial is multiplication by
$\epsilon^u-1$.  Write $u=v p^j$, where $v\in\Z$ is prime to $p$.
For $u\ne0$, characteristic $p$ and $|v|_p=1$ give
\[
 |\epsilon^u-1|=\theta^{p^j},\qquad p^j\leq |u|_{\mathbb R}.
\]
For $u>0$, the outer-radius estimate is
\[
 \frac{(b'/b)^u}{|\epsilon^u-1|}
 \leq\theta^{u-p^j}\leq1.
\]
For $u<0$, the inner-radius estimate is
\[
 \frac{(r'/r)^u}{|\epsilon^u-1|}
 \leq\theta^{|u|-p^j}\leq1.
\]
Thus division by $\epsilon^u-1$ after restriction defines a
continuous contracting homotopy on all nonzero monomials.  The
$c_0$ condition is preserved by these inequalities.  The zero
monomial has zero differential, giving the two stated terms;
the second is $\Hom(\Z_p(1),R)$ and hence $R(-1)$.

The two-term model follows by applying continuous Hom to
\[
0\to\F_p[[\Z_p]]\xrightarrow{\gamma-1}\F_p[[\Z_p]]\to\F_p\to0.
\]
Its comparison with the completed bar resolution is continuous.
Affinoid perfectoid acyclicity
\cite[Proposition~8.8]{ScholzeDiamonds} then identifies it with the
cohomology of the annulus.  The displayed coefficientwise
inequalities are unchanged for the supremum norm on
$C_{\hthreects}(T,R)$, proving the parameter assertion.
\end{proof}

\begin{lemma}[The analytic inverse limits]\label{h3:ana:analytic-limits}
Let $M_0\leftarrow M_1\leftarrow\cdots$ be a tower of complete
nonarchimedean Banach spaces with dense transition maps of norm at
most one.  Its derived inverse limit has no higher cohomology.
The assertion remains true after applying $C_{\hthreects}(T,-)$.
In particular it applies to the tilted Laurent algebras on an
exhaustion by closed smaller annuli or polyannuli.
\end{lemma}
\begin{proof}
It suffices to show that the Roos difference map
\[
 \prod_n M_n\longrightarrow\prod_n M_n,\qquad
 (x_n)\longmapsto(x_n-f_nx_{n+1})
\]
is surjective.  Given $(y_n)$, construct a solution of the first
$j$ equations at stage $j$.  To add the next equation, choose
$x_{j+1}$ so that $f_jx_{j+1}$ approximates $x_j-y_j$ to within
$2^{-j}$.  Set the new $x_j$ equal to $y_j+f_jx_{j+1}$ and
propagate this correction backwards through the earlier equations.
The transition norm bound makes the change in each earlier
coordinate at most $2^{-j}$.  Every fixed coordinate therefore
converges, and the resulting tuple solves all the equations.
This proves the claim.

Density persists for continuous functions on $T$: approximate
on a finite clopen partition and lift approximations to the
finitely many values.  Completeness and the norm bound persist
as well, so the proof applies unchanged.  Fractional Laurent
polynomials belong to every algebra in the specified exhaustion
and are dense in each of them, verifying its hypotheses.
\end{proof}

\begin{lemma}[Affine space and zero-section purity]\label{h3:ana:affine-purity}
For every $d\geq1$ and every $S,T$ as above, the natural maps give
\[
 R\Gamma(\mathbb A^d_S,\mathcal O^\flat)=R,
 \qquad
 R\Gamma_{0}(\mathbb A^d_S,\mathcal O^\flat)
   =R(-d)[-2d].
\]
The first identification is the map of constants.  The second is
the product of the normal Kummer classes.  It is natural under
translation of the section and invertible linear changes of the
normal coordinates.  The equalities are relative equalities,
including the continuous parameters $T$.
\end{lemma}
\begin{proof}
We first treat a line.  An untilted disc is covered by the
perfectoid disc obtained by adjoining all roots of its coordinate.
A section of $\mathcal O^\flat$ on the original disc pulls back to
an invariant section on this cover.  The monomial computation in
\Cref{h3:ana:annulus-contraction}, now with nonnegative exponents,
shows that every such section is constant.  Consequently the
degree-zero cohomology of a relative disc is $R$.

Restrictions from sufficiently large discs to a fixed smaller disc
kill positive cohomology.  To see this without making an assertion
about the branched cover at the center, choose a point $a$ outside
the smaller disc and inside the larger one.  The inclusion of the
smaller disc factors through two nested annuli centered at $a$,
with the radius gaps of \Cref{h3:ana:annulus-contraction}.  The larger
disc can be chosen to contain these annuli.  Restriction between
them kills the nonconstant part and all degrees above one.
The remaining degree-one class is the mod-$p$ Kummer class of
$z-a$.  On a sufficiently small disc about zero it vanishes:
if $|z/a|<|p|^{p/(p-1)}$, the binomial series for
$(1-z/a)^{1/p}$ converges, and a $p$-th root of $-a$ supplies
a first $p$-th root of $z-a$.  All choices can be made using
constants from $C$, with bounds independent of $S$ and $T$.

Exhaust the line by such discs.  For each positive degree the
inverse system of cohomology groups is pro-zero, and degree zero
is the constant system $R$.  The countable derived-limit sequence
therefore proves the first assertion for a line.  In particular,
no assertion that density alone kills a derived limit is used.

Exhaust the punctured line by annuli whose inner radii tend to zero
and outer radii tend to infinity, taking a cofinal subsequence with
the stated radius gaps.  The homotopies of
\Cref{h3:ana:annulus-contraction} show that its cohomology is
$R\oplus R(-1)[-1]$.  The constants identify the first summand.
The fiber of restriction from the line to the punctured line is
therefore $R(-1)[-2]$.  Its generator is the boundary of the
Kummer class.  Multiplication of a coordinate by an invertible
constant changes its Kummer class by the class of that constant,
which vanishes locally after adjoining a root.  Translation simply
changes the chosen section.  This proves the asserted naturality.

For $d$ coordinates apply the relative line calculation one
coordinate at a time.  Successive cohomology with support in their
zero loci gives the product of the $d$ Kummer boundary classes,
in degree $2d$ with twist $-d$.  Invertible changes of coordinates
preserve this local fundamental class: elementary additions are
checked on the complement by the same nearby-coordinate Kummer
calculation, diagonal changes were just checked, and coordinate
permutations permute classes of even degree.  These operations
generate the invertible linear changes locally on the base.
The constructions and the limiting arguments above are uniform
on continuous profinite families, proving the relative statement.
\end{proof}

\subsection{Tubes, distributions, and translation orientation}

The next support calculation replaces the single zero section by a
locally profinite family of sections.  We surround the family by disjoint
tubes and refine those tubes to the support.  The resulting support
classes will be identified with compatible residue assignments: when
tubes merge, their residues add.  The distribution module below records
this summation rule.
For a compact totally disconnected space $K$ define
\[
 \mathcal D(K,R)=
 \varprojlim_{\mathscr P}\bigoplus_{U\in\mathscr P}R,
 \qquad
 (a_V)_{V\in\mathscr Q}\longmapsto
 \left(\sum_{V\subset U}a_V\right)_{U\in\mathscr P},
 \tag{$\ast$}\label{h3:ana:distribution-definition}
\]
where $\mathscr Q$ refines the finite clopen partition
$\mathscr P$.  For a locally compact totally disconnected space
$Z=\coprod_i K_i$, put
$\mathcal D_{\mathrm{lf}}(Z,R)=\prod_i\mathcal D(K_i,R)$.
Equivalently this is the continuous dual, in this product sense,
of compactly supported locally constant functions on $Z$.
This description makes it independent of the decomposition.
It imposes no common norm bound on the different components.
The transition maps in \eqref{h3:ana:distribution-definition} are
split surjections, by choosing one child of each partition member.
For a profinite $T$ there are natural identifications
\begin{equation}
 C_{\hthreects}(T,\mathcal D_{\mathrm{lf}}(Z,R))
 =\mathcal D_{\mathrm{lf}}(Z,C_{\hthreects}(T,R)).
 \label{h3:ana:distribution-parameters}
\end{equation}
They follow directly by commuting continuous maps with products
and inverse limits.

\begin{lemma}[Support in a family of rational lattices]\label{h3:ana:tube-purity}
The following supports in an untilted relative affine line have
cohomology
\[
 R\Gamma_Z(\mathbb A^1_S,\mathcal O^\flat)
   \simeq\mathcal D_{\mathrm{lf}}(Z,R)(-1)[-2]:
\]
\begin{enumerate}[label=\textup{(\roman*)}]
\item a constant finite extension $F\subset C$ of $\Q_p$;
\item the image of $(a,b)\mapsto az+b$, with $a,b\in\Q_p$
and $z$ a section of the relative affine line whose value in each
geometric untilt lies outside $\Q_p$; one may also restrict $a$ to a
compact open subset.
\end{enumerate}
These are relative identifications, natural for the translations
and linear changes preserving the indicated support.  They retain
all profinite parameters.  Taking products of the first support
in two separate coordinates gives
\[
 R\Gamma_{F^2}(\mathbb A^2_S,\mathcal O^\flat)
   \simeq\mathcal D_{\mathrm{lf}}(F^2,R)(-2)[-4].
\]
\end{lemma}
\begin{proof}
We describe both the neighborhoods and the limit that computes
their support cohomology.  On a geometric fiber the map in (ii)
is a continuous linear injection from the finite-dimensional
$\Q_p$-Banach space $\Q_p^2$ into the geometric untilt.  Its induced
norm is equivalent to the maximum norm, and hence is bounded above
and below on the unit sphere.  The lower bound can be made uniform on a
neighborhood of that fiber.  Indeed the parameter unit sphere
in $\Q_p^2$ is compact; for each of its points the corresponding
$az+b$ is nonzero, and finitely many neighborhoods give a common
positive lower bound.  Shrink the base by their intersection.
Homothety by powers of $p$ now gives constants $c_-,c_+>0$ such that
\[
 c_-\max(|a|,|b|)\leq |az+b|
       \leq c_+\max(|a|,|b|)
\]
uniformly on that neighborhood.  The same assertion for (i) is
the norm comparison for the fixed embedding $F\subset C$.

Write $V$ for the parameter space and $d=\dim_{\Q_p}V$;
thus $d=[F:\Q_p]$ in \textup{(i)} and $d=2$ in \textup{(ii)}.
For the compact open restriction on $a$, use $V=\Q_p^2$
and start at a level at which the restricted strip is a
union of lattice cosets; compact openness permits this.
Choose a lattice adapted to the norm, as follows.  At the
chosen geometric point, the norm of the injection is locally
constant and nowhere zero on the compact parameter unit sphere,
so it takes only finitely many values there.  Choose a radius
in a gap and let $\Lambda$ be its closed parameter ball.
This is a $\Z_p$-lattice.  After a further rational restriction
of the base there are constants $\rho_-<\rho<\rho_+$ such that
\[
 |\ell(v)|\leq\rho_-\quad(v\in\Lambda),\qquad
 |\ell(v)|\geq\rho_+\quad(v\notin\Lambda),
\]
where $\ell$ denotes the indicated injection.  Only finitely
many inequalities are needed for this restriction: use a
$\Z_p$-basis of $\Lambda$ for the upper bound and representatives
$w_j$ of the nonzero cosets of $(p^{-1}\Lambda)/\Lambda$
for the lower bound.  Indeed, for $v\notin\Lambda$, a
nonnegative power of $p$ takes $v$ to
$v'=w_j+\lambda\in p^{-1}\Lambda\setminus\Lambda$.
The strict gap gives
$|\ell(v')|=|\ell(w_j)|\geq\rho_+$, and scalar rescaling
gives the desired bound for $v$.  These are inequalities
on the rational base neighborhood and hence hold on
every perfectoid test over it.

For every $n$ take discs of radius $|p|^n\rho$ about the
images of representatives modulo $p^n\Lambda$.
Distinct centers are separated by at least $|p|^n\rho_+$,
so the discs are disjoint and nested under refinement.
Only finitely many meet any bounded relative affine disc:
the lower norm bound restricts their parameters to a bounded,
hence compact, subset of $V$, which has finitely many cosets
modulo $p^n\Lambda$.  Thus the family is locally finite.

We verify its limiting support on an arbitrary affinoid
perfectoid test $S'\to S$.  Choose two buffered radii between
$\rho_-$ and $\rho_+$.  A section $w$ belonging to every
smaller closed tube system belongs to every larger open
tube system.  The inverse images of the latter's disjoint
components are open and closed in $S'$.  They define
compatible locally constant maps
$c_n:S'\to V/p^n\Lambda$.  Quasicompactness gives finitely
many values at each level.  Since
$V=\varprojlim_n V/p^n\Lambda$ with its locally profinite
topology, these maps give a continuous parameter $c:S'\to V$.
Choose representatives of the $c_n$.  Their images under
$\ell$ converge uniformly, since the successive differences
have norm at most $|p|^n\rho_-$.  The tube condition gives
$\|w-\ell(c_n)\|\leq |p|^n\rho_{\mathrm{outer}}\to0$;
completeness therefore gives $w=\ell(c)$.  Conversely a
continuous parameter $c$ belongs to all the smaller tube
systems by the upper bound on $\Lambda$.  The buffered
neighborhoods consequently have the required intersection
as relative v-sheaves, including analytic boundary points.
Their complements exhaust the desired complement, and
support cohomology is their derived inverse limit.

The exterior of a disc is exhausted by annuli.  By
\Cref{h3:ana:analytic-limits} its cohomology is the two-term
Kummer complex of Laurent series converging on that exterior.
Its degree-zero invariants are $R$.  Taking the fiber from
the affine line consequently puts the tube support entirely
in degree two.  Constant-monomial extraction gives a split
surjective residue map in that degree to $R(-1)$.
Its kernel is the nonconstant cokernel of the difference
operator.

There is a canonical description of this residue: map to
the direct limit of support cohomology over all larger
enclosing discs.  The annulus homotopy identifies that
limit with $R(-1)$.  Recentring gives the same identification,
since sufficiently far out the ratio of the two coordinates
has a first $p$-th root.  Inclusions of discs therefore
commute with residues to this common target.  Excision
identifies disjoint children with a direct sum, so the
residue of their merger is literally the sum of their
residues.  Surjectivity is also witnessed by the class of
any section inside the tube, which has residue one.

The kernels of the residue maps form a pro-zero system with
a fixed level gap.  Choose $j$ such that $|p|^j<\theta$, with
$\theta$ as in \Cref{h3:ana:annulus-contraction}; increase $j$
if necessary for the two buffered radii.  The ratio of a
level-$(n+j)$ child's radius to its level-$n$ parent's radius
is $|p|^j$, independently of $n$ and of the parent coset.
Recenter the parent disc at this child's center.  Its exterior
is unchanged because the center difference is strictly smaller
than the parent's radius.  On restriction between these
exteriors the inverse of the difference operator is bounded
on negative monomials by the fixed radius gap.  For positive
monomials use a source outer radius at least the target outer
radius divided by $\theta$, which is always available on an
exterior.  Thus the entire child residue kernel maps to zero
in the parent's support cohomology.

Each parent has exactly $p^{dj}$ children.  Finite summation
shows that the map from all level-$(n+j)$ residue kernels to
all level-$n$ kernels is zero.  For infinitely many parents
the homotopies operate coordinatewise in the product topology;
they require no common bound on the coefficient values.
For a profinite $T$ the same inequalities apply in each
$C_{\hthreects}(T,R)$ supremum norm.  Thus one level gap annihilates
the kernels on every profinite test, and not merely on
geometric fibers.

This proves that the actual system of tube support cohomology
is pro-isomorphic to the system of finite sums $R(-1)$ under
summation.  It does not require that the classes of different
sections in a fixed tube be equal.  Those classes can differ
by a residue-zero class before passage to the limit.
The residue system has split surjective transitions, while
the discarded systems are pro-zero.  Both their $\varprojlim^1$
terms consequently vanish.  Locally finite disjoint components
give products by excision.  This proves the formula
\eqref{h3:ana:distribution-definition} with its product topology.

All norm estimates and all residue maps were made over the
base neighborhood.  Applying them to $C_{\hthreects}(T,R)$ proves
the parameter assertion, and descent glues the neighborhoods
of the base.  The product statement follows by performing the
two relative support calculations successively; its partition
system consists of rectangles and is cofinal in the finite
clopen partitions of the product.
\end{proof}

\begin{lemma}[Translation cohomology]\label{h3:ana:translations}
Let $F/\Q_p$ have degree $m$, acting by translations.  Then
\[
 R\Gamma_{\hthreects}(F,R)=R,\qquad
 R\Gamma_{\hthreects}(F,\mathcal D_{\mathrm{lf}}(F,R))
       =R\otimes_{\F_p}\operatorname{or}_F[-m].
\]
Here $\operatorname{or}_F$ is the inverse orientation of the
$m$-dimensional translation group.  With the convention that
$a$ acts on points by $z\mapsto az$, a scalar
$a\in\Z_p^\times$ acts on this line by $\bar a^{-m}$.
The assertions and their support maps hold with continuous
profinite parameters.
\end{lemma}
\begin{proof}
Choose a lattice $K\simeq\Z_p^m$ in $F$.  On its distributions
the finite Koszul complex is the Koszul cochain complex for
\[
 R[[X_1,\ldots,X_m]],\qquad X_i=\gamma_i-1.
\]
The power-series notation here means the inverse limit of finite
group algebras, as in \eqref{h3:ana:distribution-definition}, hence
a product of coefficient copies.  Multiplication by a variable
is injective and its quotient is obtained by setting that
variable to zero.  Successively applying these statements
shows that the Koszul cochain complex has just its top
cohomology, the augmentation quotient $R$, in degree $m$.
The coefficient splitting by powers of $X_i$ makes these
statements exact also with continuous parameters.

There is an identification of translation modules
\[
 \mathcal D_{\mathrm{lf}}(F,R)
  =\operatorname{Coind}_{K}^{F}\mathcal D(K,R).
\]
Indeed $F/K$ is discrete, and both sides are the product of
one copy of $\mathcal D(K,R)$ for each coset, with the same
translation rule.  Evaluation on $K$, and the bar contraction
using representatives for these open cosets, give continuous
cohomological Shapiro for this coinduced module.  This proves
the second formula.  Changing a lattice basis acts on the
top cochain generator by the inverse determinant of its
change matrix.  Scalar multiplication by $a$ has determinant
$a^m$, proving the asserted character.

For the first formula, exhaust $F$ by $p^{-j}K$.  In every cochain
degree, restriction between these open-and-closed subgroups is
surjective by extension by zero.  The inverse limit is the continuous
cochain complex on $F$, since these open subgroups cover $F$; the same
statements hold with the parameter $T$.  On trivial
characteristic-$p$ coefficients the Koszul differentials
are zero, and compact-lattice cohomology is the exterior
algebra on the continuous linear forms of the lattice.
Restriction from $p^{-j-1}K$ to $p^{-j}K$ multiplies each
degree-one form by $p$ and is therefore zero in every positive
degree.  In degree zero restriction is the identity.
The derived inverse limit is consequently $R$.
The same finite complexes and the same vanishing transition
maps compute the parameter versions.
\end{proof}

\begin{proposition}[The punctured negative space]\label{h3:prop:punctured-cohomology}
For $m=1,2$ the natural constants map fits into a fiber sequence
\[
 R\longrightarrow R\Gamma(N_m^*,\mathcal O^\flat)
   \longrightarrow
 R(-1)\otimes\operatorname{or}_F[-m-1].
\]
It is equivariant for scalar units and natural in $S$ and in
profinite parameters.  Thus the only cohomology groups are
$R$ in degree zero and an orientation line in degree $m+1$.
The scalar $a\in\Z_p^\times$ has weight $\bar a^{-m}$ on
the latter.  Reversing the action convention reverses all
these exponents simultaneously.
\end{proposition}
\begin{proof}
Use the support fiber sequence for
$\underline F\subset\mathbb A^1_C$ in
\Cref{h3:ana:relative-bc}.  Its middle term is $R$ by
\Cref{h3:ana:affine-purity}, and its first term is
$\mathcal D_{\mathrm{lf}}(F,R)(-1)[-2]$ by
\Cref{h3:ana:tube-purity}.  Take the translation cohomology
of \Cref{h3:ana:translations}.  The resulting support term
has degree $m+2$; rotation gives the stated cofiber in
degree $m+1$.  The map from $R$ throughout this construction
is the coefficient unit.  Scalar multiplication preserves
the normal Kummer class and contributes exactly the
translation orientation computed in that lemma.
\end{proof}

\subsection{Independent pairs and compact frame cohomology}

The punctured calculation has produced the two cohomology rows for
coheight one.  For the ordinary stratum, we must instead remove the
dependent pairs in $N_1^2$.  We compute the two support terms with
their frame characters before taking the compact frame-group
cohomology.

Put $N=N_1$, $E=E_1$, and $G_f=\GL_2(\Z_p)\times P_1$.
Let $D\subset N^2\setminus\{0\}$ be the locus of proportional
pairs, so that $E=(N^2\setminus\{0\})\setminus D$.
The direction of a proportional pair belongs to
$\mathbb P^1(\Q_p)$.

\begin{lemma}[The two support terms]\label{h3:ana:ordinary-support}
The support at the double origin in $N^2$ is a line in degree
six, of scalar weight $-2$.
The support in $D$ has cohomology in degrees three and five.
In each degree it is a distribution module on
$\mathbb P^1(\Q_p)$ with a one-dimensional continuous fiber.
The scalar weights of these two fibers are respectively
$-1$ and $-2$.

These descriptions are equivariant relative descriptions,
with continuous profinite parameters.  For the stabilizer
of the direction $(1,0)$, written
\[
 B=\left\{\begin{pmatrix}a&b\\0&d\end{pmatrix}
                \in\GL_2(\Z_p)\right\},
\]
the two fibers have characters $\bar d^{-1}$ and
$\bar a^{-1}\bar d^{-1}$.  An element $v\in P_1=\Z_p^\times$
has characters $\bar v$ and $\bar v^2$, respectively.
Normal Kummer twists are understood in these descriptions.
\end{lemma}
\begin{proof}
The preimage of the double origin under
$\mathbb A^2_C\to N^2$ is $\Q_p^2$.
The product support calculation gives its distributions in
degree four.  Taking the two translation directions adds
degree two and inverse determinant orientation, giving
degree six and scalar weight $-2$.

On a slope chart, lift the first coordinate to
$z_1\notin\Q_p$.  The support in the second coordinate is
\[
 z_2=a z_1+b,\qquad a,b\in\Q_p,
\]
with $a$ in a compact open slope patch when desired.
\Cref{h3:ana:tube-purity} gives distributions in $(a,b)$ in
normal support degree two.  Taking translations in the
second coordinate is the $b$-translation calculation of
\Cref{h3:ana:translations}.  It leaves distributions in $a$
and adds degree one with inverse orientation.  First-coordinate
translation replaces $b$ by $b-ac$ and leaves $a$ fixed.
It consequently preserves the resulting orientation line.
The remaining base is $N^*$; \Cref{h3:prop:punctured-cohomology}
adds its degrees zero and two.  We obtain degrees three and
five, and weights $-1$ and $-2$ with the same sign.
To justify this computation with distributions in the slope,
first use a finite clopen slope partition.  It gives a finite
sum of the relative punctured calculation.  Then take the
partition limit.  The summation maps are split and the
discarded residue kernels are uniformly pro-zero by
\Cref{h3:ana:tube-purity}; hence this limit introduces no
additional cohomology.  This also specifies the topology
of the resulting distribution rows.

For equivariance, a point mass at a direction is the support
map from that direction's tubes.  These maps commute with
linear changes of the pair.  Finite sums of point masses
are dense in the finite-partition topology: on any finite
partition every prescribed collection of values is realized
by such a sum.  Transport of the point masses therefore
determines transport on the full distribution module.
At the direction $(1,0)$ the tangential coordinate is changed
by $a$ and the normal coordinate by $d$.  The inverse
translation orientations just computed are $\bar d^{-1}$
and $\bar a^{-1}\bar d^{-1}$; the upper-right entry acts
trivially on both fibers.  The quotient frame $v$ scales
both coordinates by $v^{-1}$, giving the stated $P_1$
characters.  These are finite characters.  The calculations
of tubes and translations were relative, so this transport
argument also applies with all the indicated parameters.
\end{proof}

\begin{proposition}[The ordinary constants map]\label{h3:prop:ordinary-constants}
For $p\geq5$, the structure map $Y_1\to BG_f$ and the
coefficient unit induce an equivalence
\[
 R\Gamma_{\hthreects}(G_f,R)
   \xrightarrow{\ \simeq\ }
 R\Gamma(Y_1,\mathcal O^\flat).
\]
This is an equivalence of derived algebras and is natural
over $S$, with continuous profinite parameters.  It retains
both frame-group actions before taking their cohomology.
In particular the degree-three class from the rank-two
Tate frame is carried by this actual map from constants.
\end{proposition}
\begin{proof}
By \Cref{h3:ana:affine-purity,h3:ana:translations}, constants
compute the cohomology of $N^2$.  Remove the double origin
and then the proportional locus.  These two support fiber
sequences are equivariant for $G_f$ and their maps from
constants are the coefficient unit.

The central subgroup $\mu_{p-1}I_2\subset\GL_2(\Z_p)$
acts on their support cohomology by the characters of
weights $-1$ and $-2$ in \Cref{h3:ana:ordinary-support}.
Neither character is trivial when $p\geq5$.  Its averaging
idempotent is defined in characteristic $p$, so its derived
invariants are exact and annihilate both support complexes.
Taking the remaining frame-group cohomology therefore
turns the constants arrow into an equivalence.
The arrow was induced by the structure map and coefficient
unit, which are multiplicative.  Its equivalence on
underlying complexes proves the assertion for derived
algebras as well.
\end{proof}

For later finiteness arguments we need twists even when
the preceding averaging no longer kills the support.
The appropriate induction is induction of measures.

\begin{lemma}[Cohomology of induced measure modules]\label{h3:ana:measure-induction}
Let $G$ be a compact $p$-adic analytic group without
$p$-torsion, let $B\subset G$ be a closed analytic subgroup
without $p$-torsion, and suppose $G\to G/B$ has continuous
local sections.  If $V$ is a finite-dimensional continuous
$k$-representation of $B$, then
\[
 M=k[[G]]\widehat\otimes_{k[[B]]}V
\]
has finite-dimensional continuous $G$-cohomology in every
degree.  The same cohomology after replacing coefficient
copies of $k$ by $R$ is the corresponding finite-dimensional
cohomology tensored with $R$, including continuous
profinite parameters.
\end{lemma}
\begin{proof}
Write $g=\dim G$.  Analytic group duality identifies
cohomology with homology in complementary degree:
\[
 H^a_{\hthreects}(G,M)
   \simeq H^{\hthreects}_{g-a}
        (G,\operatorname{or}_G\otimes_k M).
\]
The orientation module is retained in this formula.
One obtains it by dualizing a finite projective completed
group-ring resolution of $k$ and reversing its degrees.
Such resolutions exist for the stated groups by
\cite[Proposition~3.3]{AW}; their dualizing module is the top inverse
adjoint orientation
\cite[Proposition~4.16, Remark~4.23, and Proposition~4.40]{BGHS}.
Thus the formula applies to compact
measure modules, not just to finite coefficients.

Completed homological Shapiro gives
\[
 H^{\hthreects}_{g-a}(G,\operatorname{or}_G\otimes M)
   \simeq H^{\hthreects}_{g-a}
       (B,\operatorname{or}_G|_B\otimes V).
\]
For completeness, local sections identify $k[[G]]$, as a
right completed $k[[B]]$-module, with a completed free
module on $G/B$.  Inducing a completed projective
$B$-resolution and then taking $G$-coinvariants is
therefore the same complex as taking $B$-coinvariants.
This proves the displayed Shapiro identification.
The last homology is finite-dimensional, since a finite
projective $B$-resolution has finite-dimensional terms
after tensoring with the finite-dimensional fiber.
This argument includes the duality shift; it makes no
unshifted cohomological Shapiro assertion for $M$.

To justify extension to $R$, start with compact $k$-modules
and apply the exact functor
\[
 A\longmapsto\Hom_k(A^\vee,R),\qquad
 A^\vee=\Hom_{\hthreects,k}(A,k).
\]
After a choice of a basis of $A^\vee$ its value is a
product of copies of $R$.  Exactness follows by splitting
sequences of the discrete vector spaces $A^\vee$.
Finite projective resolutions use only finite powers
and idempotent summands, so their comparison, duality,
and Shapiro maps commute with this functor.  Apply the
same argument to $C_{\hthreects}(T,R)$, or use
\eqref{h3:ana:distribution-parameters}.
\end{proof}

\begin{corollary}[Twisted ordinary finiteness]\label{h3:ana:twisted-finiteness}
If a coefficient line on $Y_1$ becomes a finite-character
constant line on $E_1$, its geometric cohomology is
finite-dimensional in every degree and vanishes outside
a finite range.  This applies to all powers of the
periodicity line and all finite tame determinant twists.
For a profinite parameter $T$, the corresponding cohomology
is $C_{\hthreects}(T,-)$ of this finite-dimensional answer.
\end{corollary}
\begin{proof}
The ambient constants and the origin-support line have
finite-dimensional fibers with finite-character action.
The two remaining support rows in
\Cref{h3:ana:ordinary-support} are completed measure induction
from $B\times P_1$ to $G_f$.  To check the local-section
hypothesis, a primitive vector over $\Z_p$ can, on either
standard slope chart, be completed continuously to an
invertible matrix.  These sections also show that
$\GL_2(\Z_p)$ acts transitively on $\mathbb P^1(\Q_p)$.
Its stabilizer is the displayed upper triangular group.
Both $G_f$ and $B\times P_1$ are compact analytic without
$p$-torsion when $p\geq5$: a matrix of order $p$ in
dimension two would require the degree-$p-1$ cyclotomic
polynomial, and $\Z_p^\times$ has no $p$-torsion.
\Cref{h3:ana:measure-induction} proves finiteness of their
support cohomology.  The two support sequences prove the
claim for the independent locus.

On the connected Honda marking the periodicity line is
framed by the leading derivative of the formal-group
isomorphism.  At height one its change of frame is
the reduction $\Z_p^\times\to\F_p^\times$.  The determinant
identity of \Cref{h3:prop:integral-frames} expresses a tame
middle determinant character through the two frame
characters.  The indicated coefficient lines therefore
have exactly the finite-character property used above.
\end{proof}

\subsection{Descent to the finite constant field}

The geometric calculations are over $C^\flat$.  We now return to the
finite field $k$ while retaining the stack maps and the residual group
actions.  This step uses base-field Frobenius on the $k$-defined
spaces; the chosen untilt chart remains a device for computing their
geometric cohomology.

For a $k$-defined space $X$, base-field Frobenius means the
map $\hthreeid_X\times\varphi_q$ on
$X\times_{\operatorname{Spd}k}\operatorname{Spd}C^\flat$.
In a coordinate chart defined over $k$ it raises coefficients
to the $q$-th power and fixes the variables.  It is distinct
from Frobenius of all the variables in a perfected coefficient
ring.  A chosen untilt-divisor presentation need not be
preserved by this map.

\begin{lemma}[Base-field Frobenius with parameters]\label{h3:ana:base-frobenius}
Let $q=|k|$, and let $M_k$ be the analytic function space over $k$
on a chosen perfected punctured ball or on the perfected ordinary
polyannular open.  Let $M_C$ be the corresponding function space after
extension from $k$ to $C^\flat$.  Let $\varphi$ act as the $q$-th power
on coefficients and fix the variables.  Then
\[
 0\longrightarrow M_k\longrightarrow M_C
      \xrightarrow{\varphi-1}M_C\longrightarrow0
\]
is exact, including after $C_{\hthreects}(T,-)$ for a profinite
space $T$.  The same statement applies to finite products
of the punctured-ball charts at finite marking levels.
Consequently the coefficient comparison is an equality
with derived base-Frobenius invariants, naturally for all
group actions defined over $k$.
\end{lemma}
\begin{proof}
Use the expansions in fractional monomials on exhausted
closed smaller polyannuli.  An invariant expansion has
each coefficient in $\F_q=k$, and hence is precisely an
element of $M_k$.  For a coefficient $a\in C^\flat$, choose
a root $b$ of $b^q-b=a$ of least absolute value.  The
Newton polygon, or the convergent series
$b=-\sum_{j\geq0}a^{q^j}$ when $|a|<1$, gives
\[
 |b|=|a|\quad (|a|<1),\qquad
 |b|\leq1\quad (|a|=1),\qquad
 |b|=|a|^{1/q}\quad (|a|>1).
\]
In particular $|b|\leq|a|$.  Solving for each coefficient
preserves the $c_0$ convergence conditions in every
smaller polyannulus norm.  This proves surjectivity and
also supplies solutions arbitrarily small whenever the
input is sufficiently small in any finite collection
of these norms.
The Frobenius map itself is continuous in the full analytic
topology: for a Gauss poly-radius $\rho$,
$\|\varphi f\|_\rho=\|f\|_{\rho^{1/q}}^q$.
The slightly larger poly-radius on the right belongs to
the same open domain.

For continuous parameter families, first approximate a
given function on a finite clopen partition of $T$ and
lift the finitely many values.  The remaining error is
small.  Solve for its coefficients with the small
branch just described, refine the partition, and repeat
with errors tending to zero in successive defining
norms.  The estimates give a uniformly convergent sum
of continuous corrections; its limit is a continuous
lift.  This argument proves exactness on parameters
directly.  For the spaces defined by an exhaustion,
the continuous Laurent-polynomial approximation and
the complete norm topologies meet the hypotheses of
\Cref{h3:ana:analytic-limits}.  Thus these analytic inverse
limits have no additional derived term.

The maps $M_k\to M_C$ and $\varphi-1$ themselves are
natural for changes of variables defined over $k$.
Their auxiliary coefficient choices need not be natural.
The exact sequence therefore identifies the original
coefficient complex with the fiber of $\varphi-1$,
equivariantly and with parameters.  If a finite marking
chart has residue field $k'/k$, geometric base change
gives a product indexed by the embeddings of $k'$.
The canonical base Frobenius permutes these factors.
On a cycle of length $e=[k':k]$, eliminating successive
coordinates reduces its difference equation to
coefficient Frobenius $q^e$ minus identity.  The preceding
estimates and continuous lifting apply with $q^e$.
Thus finite products and finite constant-field descent
retain the claimed exact sequence, with the actual
permutation of factors.
\end{proof}

Write $B_{\hthreepk}$ for the algebra of functions on $\mathcal X_1$,
the completed perfection of the ordinary locus, with the inverse-limit
topology of uniform convergence on the exhausted closed polyannular
charts.  \Cref{h3:sec:ordinary} gives an equivalent description using
finite root levels of Laurent series and a countable family of defining
norms.

\begin{theorem}[Completed cohomology and its constants maps]
\label{h3:thm:completed-cohomology}
The completed strata of \Cref{h3:thm:coordinate-comparison}
have the following properties.
\begin{enumerate}[label=\textup{(\roman*)}]
\item For the ordinary completed coefficient ring
$B_{\hthreepk}$, the natural constants map is an equivalence
\[
 R\Gamma_{\hthreects}(\GL_2(\Z_p)\times P_1,k)
   \xrightarrow{\ \simeq\ }
 R\Gamma_{\hthreects}(P_3,B_{\hthreepk}).
\]
It is the map defined by $Y_1\to B(\GL_2(\Z_p)\times P_1)$,
and is compatible with geometric constant extension,
products, coefficient Frobenius, and profinite parameters.
Every periodicity twist and finite tame character twist
has finite cohomology in each degree, in a bounded range.

\item Let $L$ be the connected height-two marking algebra,
$U\subset P_2$ an open subgroup, $B=L^U$, and
$B^{\mathrm b}=\widehat{B^{\mathrm{perf}}}$.  Finite products of
fields are allowed.  Then each
$H^a_{\hthreects}(H,B^{\mathrm b})$ is finite-dimensional over $k$,
and its degree-zero group is $k$ by constants.
The same finiteness holds for lines framed on a finite
connected-marking cover, including the periodicity and
tame determinant lines used below.
On taking the filtered colimit over marking levels there
is a natural fiber sequence
\[
 k\longrightarrow
 \colim_U R\Gamma_{\hthreects}(H,(L^U)^{\mathrm b})
   \longrightarrow k_{\lambda}[-3].
\]
The first arrow is the inclusion of constants.  The
full residual norm quotient $P_3/H=\Z_p^\times$ acts
on $k_\lambda$ by $a\mapsto\bar a^{-2}$ in the
point-action convention above.  The commuting $P_2$
action is retained.  It is smooth on the displayed
finite-dimensional rows.  The statement includes
restrictions between marking levels.  With a profinite
parameter $T$, each displayed cohomology group is replaced
by its $C_{\hthreects}(T,-)$, as specified at the start of the section.
\end{enumerate}
In particular, invariants of the order-$p-1$ subgroup
of the full norm quotient kill the extra line in
\textup{(ii)}.  This subgroup is in the quotient;
no section of $P_3\to\Z_p^\times$ is required.
\end{theorem}
\begin{proof}
First compute over $C^\flat$.  For the ordinary locus,
\Cref{h3:prop:ordinary-constants} is precisely the desired
map on the quotient-stack presentation in
\Cref{h3:prop:integral-frames}.  The perfected polyannular
exhaustion has affinoid perfectoid charts.  Its coefficient
complex is the function complex $B_{\hthreepk}$, since the
successive polynomial approximation in
\Cref{h3:ana:analytic-limits} gives the required vanishing
of the derived analytic limit.  Thus this is also the
ordinary continuous $P_3$-cochain comparison with the
stated topology.

Take derived invariants of base-field Frobenius using
\Cref{h3:ana:base-frobenius}.  On the constants side the
fiber of $\varphi-1$ is $k$.  This proves the untwisted
comparison over $k$ with its actual coefficient map.
For a twist, \Cref{h3:ana:twisted-finiteness} gives a bounded
complex with finite-dimensional geometric cohomology.
Frobenius acts semilinearly and bijectively on these
groups.  On a finite-dimensional $C^\flat$-space a
bijective $q$-semilinear operator $F$ has $F-1$
surjective and a finite-dimensional fixed space over
$k$.  Indeed write $F(x)=A x^{[q]}$ with $A$ invertible.
The system $A x^{[q]}-x=b$ is equivalent to monic
equations $x_i^q=(A^{-1}(x+b))_i$.  Their pairwise
coprime leading monomials show that the coordinate
algebra is free over the $b$-coordinate algebra on
the $q^n$ monomials with each exponent less than $q$.
The derivative is $-I$, so the map is finite \'etale,
surjective, and has kernel of order $q^n$.  The kernel
is a $k$-vector space and therefore has dimension $n$.
This proves the asserted semilinear calculation.
Hence Frobenius descent
preserves the finite-dimensional conclusion.  All
these operators commute with the frame actions and
with the coefficient maps.

For \textup{(ii)}, define Frobenius before choosing an
untilt-divisor presentation.  The space $N_2$ and its
$P_2$ and residual norm actions are defined over $k$.
Consequently $\hthreeid_{N_2}\times\varphi_q$ induces a
$q$-semilinear automorphism $\Phi$ of
\[
 \mathcal K_C=R\Gamma(N_{2,C}^*,\mathcal O^\flat),\qquad
 N_{2,C}=N_2\times_{\operatorname{Spd}k}
                       \operatorname{Spd}C^\flat,
\]
commuting with both actions.  The presentation
$N_{2,C}^*\simeq[(\mathbb A_C^1\setminus F)/F]$ computes
$\mathcal K_C$.  It is not required to commute with
$\Phi$: transporting $\Phi$ through it can change the
chosen untilt divisor.  Put
$\mathcal K_k=\fib(\Phi-1:\mathcal K_C\to\mathcal K_C)$.

Write $X_B=\operatorname{Spa}(B^{\mathrm b})$ and
$M_{B,C}=\Gamma(X_B\times_k\operatorname{Spd}C^\flat,
\mathcal O^\flat)$.  The perfected punctured-ball
exhaustion and \Cref{h3:ana:analytic-limits} give no
higher cohomology on this geometric chart.  On its
$k$-defined coordinates the canonical base Frobenius
is coefficient Frobenius, with any finite residue-field
factor permutation retained.  Thus
\Cref{h3:ana:base-frobenius}, applied in every bar degree
with parameter $H^a$, identifies
$R\Gamma_{\hthreects}(H,B^{\mathrm b})$ with the fiber of
$\Phi-1$ on $R\Gamma_{\hthreects}(H,M_{B,C})$.

The actual marking torsor of
\Cref{h3:geom:finite-level-stack} gives a $k$-defined
equivalence $[X_B/H]\simeq[N_2^*/U]$.
Base-changing this equivalence supplies the vertical
maps in the commutative diagram
\begin{equation}\label{h3:ana:frobenius-square}
\begin{tikzcd}[column sep=large]
 R\Gamma_{\hthreects}(H,M_{B,C})
   \arrow[r,"\Phi-1"]\arrow[d,"\simeq"']
 &R\Gamma_{\hthreects}(H,M_{B,C})\arrow[d,"\simeq"]\\
 R\Gamma_{\hthreects}(U,\mathcal K_C)\arrow[r,"\Phi-1"]
 &R\Gamma_{\hthreects}(U,\mathcal K_C).
\end{tikzcd}
\end{equation}
These maps come from the quotient stacks over $k$.
The diagram is natural for $P_2$-transport between
marking levels and is equivariant for the full
residual norm and the coherent Frobenius actions.
Taking fibers, and commuting the two derived limits,
gives the natural identity
\[
 R\Gamma_{\hthreects}(H,B^{\mathrm b})
      \simeq R\Gamma_{\hthreects}(U,\mathcal K_k).
\]

There are no additional rows in this fiber.
\Cref{h3:prop:punctured-cohomology} gives geometric
cohomology only in degrees zero and three.  In every
degree the fiber sequence gives
\[
 0\longrightarrow
 \operatorname{coker}(\Phi-1\mid H^{i-1}\mathcal K_C)
 \longrightarrow H^i\mathcal K_k
 \longrightarrow
 \ker(\Phi-1\mid H^i\mathcal K_C)
 \longrightarrow0.
\]
The semilinear calculation above makes the two
potentially nonzero cokernels vanish.  Thus
$\mathcal K_k$ has the two stated $k$-rows, with
degree zero equal to $k$ and trivial $P_2$ action.
Smoothness of the upper action can also be seen
directly: if $e$ is a basis, $\Phi(e)=Ae$, and
$u(e)=c(u)e$, commutation gives
$c(u)^q=c(u)$.  Hence $c(u)\in k^\times$.
The scalar norm character $\bar a^{-2}$ survives
this descent since it already takes values in
$\F_p^\times$.

A compact analytic $U\subset P_2$ has finite
completed resolutions: it has no $p$-torsion,
because an order-$p$ element in its degree-two
division algebra would require a subfield of
degree $p-1>2$.  The quotient spectral sequence
with the two finite $k$-rows therefore proves
finite-stage finiteness and identifies degree
zero with $k$.  The resolutions consist of
finite powers and idempotent summands.  The
continuous Artin--Schreier lifting above and
these finite resolutions identify the actual
parameter cohomology with $C_{\hthreects}(T,-)$ of
the point-valued finite groups.
For a line framed on a finite connected-marking
cover, the same argument has a finite-character
$U$-fiber.  The tautological connected frame
trivializes the periodicity line under $H$, and a
tame determinant character is trivial on $H$.
This proves the additional finite-stage assertion.

Finally choose a cofinal sequence of open uniform
subgroups $U$, small enough to act trivially on
these two finite rows.  Restriction to a sufficiently
deep $p$-power subgroup is zero in positive
cohomological degrees: on the mod-$p$ exterior
cohomology of a uniform group it is zero in degree
one and hence in every positive degree.  Filtered
colimits of the quotient spectral sequences
therefore retain precisely the two rows themselves.
The constants arrow gives the stated fiber
sequence.  Its scalar character was computed on
$N_2^*$ in \Cref{h3:prop:punctured-cohomology}.
By determinant normalization in
\Cref{h3:prop:integral-frames}, this scalar is the
\emph{full} reduced norm of an element of $P_3$.
The induced quotient action on $H$-cohomology is
well defined because inner conjugation by $H$
acts trivially there.  Exact averaging in the
prime-to-$p$ norm quotient gives the final assertion.
\end{proof}

\begin{remark}\label{h3:ana:witt-naturality}
The structure map used in
\Cref{h3:prop:ordinary-constants,h3:thm:completed-cohomology}
also defines
\[
 R\Gamma(G_f,W_\ell(k))\longrightarrow
 R\Gamma(Y_1,W_\ell\mathcal O^\flat)
\]
for every $\ell$.  These maps commute with Verschiebung,
Frobenius, truncation, and the maps from
$\Z/p^\ell$-constants.  Their equivalence follows
by induction on the finite Verschiebung filtration
from the length-one equivalence.  Thus the constants
comparison retains the actual frame classes through
the finite Witt comparisons used later.
\end{remark}
